% ============================================================
%  Cours : Analyse des Données
%  Public : Master GAMMA · Master BA · Licence LMAD
%  Style  : Beamer Metropolis
% ============================================================
\documentclass[10pt, aspectratio=169]{beamer}

% ─── Thème ──────────────────────────────────────────────────
\usetheme{metropolis}
\usepackage{appendixnumberbeamer}

% ─── Encodage et langue ─────────────────────────────────────
\usepackage[utf8]{inputenc}
\usepackage[T1]{fontenc}
\usepackage[french]{babel}

% ─── Mathématiques ──────────────────────────────────────────
\usepackage{amsmath, amssymb, amsthm, bm, mathtools}
\usepackage{bbm}
\usepackage{dsfont}

% ─── Tableaux ───────────────────────────────────────────────
\usepackage{booktabs}
\usepackage{array}
\usepackage{multirow}

% ─── Code R / Python ────────────────────────────────────────
\usepackage{listings}
\usepackage{algorithm}
\usepackage{algpseudocode}

% ─── Figures ────────────────────────────────────────────────
\usepackage{graphicx}
\usepackage{tikz}
\usepackage{pgfplots}
\pgfplotsset{compat=1.18}
\usetikzlibrary{
  arrows.meta, positioning, shapes.geometric,
  fit, backgrounds, calc, matrix, chains,
  decorations.pathmorphing, patterns
}

% ─── Colonnes ───────────────────────────────────────────────
\usepackage{multicol}

% ─── Couleurs ───────────────────────────────────────────────
\definecolor{deepblue}{RGB}{13,71,161}
\definecolor{deeporange}{RGB}{230,81,0}
\definecolor{darkgreen}{RGB}{27,94,32}
\definecolor{advcolor}{RGB}{180,0,0}
\definecolor{mygray}{RGB}{245,245,245}
\definecolor{thmblue}{RGB}{0,102,153}

% ─── Palette Metropolis personnalisée ───────────────────────
\setbeamercolor{frametitle}{bg=deepblue, fg=white}
\setbeamercolor{progress bar}{fg=deeporange}
\setbeamercolor{alerted text}{fg=deeporange}
\setbeamercolor{example text}{fg=darkgreen}
\setbeamertemplate{frame numbering}[fraction]
\metroset{block=fill}

% ─── Théorèmes ──────────────────────────────────────────────
% beamer-metropolis prédéfinit definition/theorem/example/proof.
% On rajoute proposition / lemma / corollary / remark en *unstarred*
% avec des noms d'environnement non conflictuels.
\theoremstyle{plain}
\newtheorem{prop}{Proposition}
\newtheorem{lem}{Lemme}
\newtheorem{coro}{Corollaire}
\theoremstyle{remark}
\newtheorem{rema}{Remarque}

% ─── Style du code R/Python ─────────────────────────────────
\lstdefinelanguage{Renh}{
  morekeywords={if,else,for,while,function,return,in,TRUE,FALSE,NULL,NA,NaN,Inf,library,require,source},
  morecomment=[l]{\#},
  morestring=[b]",
  morestring=[b]'
}
\lstset{
  language=Renh,
  backgroundcolor=\color{mygray},
  basicstyle=\ttfamily\scriptsize,
  keywordstyle=\color{deepblue}\bfseries,
  stringstyle=\color{deeporange},
  commentstyle=\color{gray}\itshape,
  showstringspaces=false,
  numbers=left,
  numberstyle=\tiny\color{gray},
  frame=single,
  framesep=3pt,
  breaklines=true,
  tabsize=2,
  inputencoding=utf8,
  extendedchars=true,
  literate={é}{{\'e}}1 {è}{{\`e}}1 {ê}{{\^e}}1 {à}{{\`a}}1
           {â}{{\^a}}1 {ç}{{\c c}}1 {ô}{{\^o}}1 {î}{{\^\i}}1
           {É}{{\'E}}1 {È}{{\`E}}1 {Ê}{{\^E}}1 {À}{{\`A}}1
           {ù}{{\`u}}1 {û}{{\^u}}1 {ï}{{\"\i}}1 {œ}{{\oe}}1
           {’}{{'}}1 {–}{{-}}1 {—}{{-}}1 {«}{{<<}}2 {»}{{>>}}2,
}

% ─── Macros ─────────────────────────────────────────────────
\newcommand{\adv}{\textcolor{advcolor}{\textsf{\textbf{[ADV]}}}\;}
\DeclareMathOperator*{\argmin}{arg\,min}
\DeclareMathOperator*{\argmax}{arg\,max}
\DeclareMathOperator{\Tr}{Tr}
\DeclareMathOperator{\Var}{Var}
\DeclareMathOperator{\Cov}{Cov}
\DeclareMathOperator{\Cor}{Cor}
\DeclareMathOperator{\rang}{rang}
\newcommand{\R}{\mathbb{R}}
\newcommand{\N}{\mathbb{N}}
\newcommand{\E}{\mathbb{E}}
\newcommand{\1}{\mathds{1}}
\newcommand{\Loss}{\mathcal{L}}

% ─── Page de titre ──────────────────────────────────────────
\title{Analyse des Données}
\subtitle{ACP · AFC · ACM · Classification · AFD}
\author{Aymen Ben Brik}
\institute{Esprit School of Business\\
           Master GAMMA · Master BA · Licence LMAD}
\date{Année académique 2025--2026}

% ============================================================
\begin{document}
% ============================================================

\maketitle

% ── Table des matières ──────────────────────────────────────
\begin{frame}{Plan du cours}
  \setbeamertemplate{section in toc}[sections numbered]
  \tableofcontents[hideallsubsections]
\end{frame}

% ── Chapitres ───────────────────────────────────────────────
% ============================================================
%  Chapitre 1 — Introduction et Statistiques Descriptives Multivariées
% ============================================================
\section{Introduction et Statistiques Descriptives Multivariées}

% ────────────────────────────────────────────────────────────
\subsection{1.1 Position de l'analyse des données}
% ────────────────────────────────────────────────────────────
\begin{frame}{1.1 — Qu'est-ce que l'analyse des données ?}
\begin{columns}[T]
\column{0.55\textwidth}
\begin{definition}
L'\textbf{analyse des données} (\textit{data analysis}) est l'ensemble des
techniques exploratoires \alert{descriptives}, fondées sur la
\textbf{décomposition matricielle} et la \textbf{géométrie des nuages de points},
visant à révéler la structure d'un tableau de données
$\mathbf{X} \in \R^{n \times p}$ \emph{sans modèle probabiliste}.
\end{definition}

\vspace{0.3cm}
\textbf{Tradition française} (Benzécri, années 1960) :
\begin{itemize}\small
  \item \alert{Géométrique} (inertie, projections orthogonales)
  \item \alert{Descriptive} (pas d'hypothèses inférentielles)
  \item \alert{Visuelle} (cartes factorielles)
\end{itemize}
\column{0.42\textwidth}
\begin{tikzpicture}[scale=0.85, font=\scriptsize,
  box/.style={draw=deepblue, fill=deepblue!10, rounded corners,
              minimum width=3.3cm, minimum height=0.65cm, align=center}]
  \node[box]                      (stat) at (0, 2.2) {Statistique\\inférentielle};
  \node[box, fill=deeporange!10, draw=deeporange]  (ad)   at (0, 0.9) {Analyse\\des données};
  \node[box, fill=darkgreen!10, draw=darkgreen]    (ml)   at (0,-0.4) {Machine\\Learning};
  \node[font=\tiny, gray, align=center] at (4.4, 2.2) {modèle\\$+$ tests};
  \node[font=\tiny, gray, align=center] at (4.4, 0.9) {description\\géométrique};
  \node[font=\tiny, gray, align=center] at (4.4,-0.4) {prédiction\\$+$ optimisation};
  \draw[->, thick, gray] (stat.south) -- (ad.north);
  \draw[->, thick, gray] (ad.south)   -- (ml.north);
\end{tikzpicture}
\end{columns}
\end{frame}

% ────────────────────────────────────────────────────────────
\begin{frame}{1.1 — Objectifs de l'analyse des données}
\begin{columns}[T]
\column{0.49\textwidth}
\textbf{Trois grands objectifs}
\begin{enumerate}
  \item \textbf{Réduire la dimension} : remplacer $p$ variables corrélées
        par $q \ll p$ \emph{composantes synthétiques}.
  \item \textbf{Représenter graphiquement} les individus et les variables
        sur des plans factoriels.
  \item \textbf{Classifier} (typer) les individus en groupes homogènes.
\end{enumerate}

\vspace{0.3cm}
\textbf{Méthodes selon le type de données}
\begin{itemize}\small
  \item Quantitatives $\to$ \alert{ACP}
  \item Tableau de contingence $\to$ \alert{AFC}
  \item Qualitatives $\to$ \alert{ACM}
  \item Mixtes $\to$ AFDM, AFM
\end{itemize}
\column{0.49\textwidth}
\begin{block}{Place de ce cours}
  Articulation entre :
  \begin{itemize}\small
    \item \textbf{algèbre linéaire} (LMAD, théorèmes)
    \item \textbf{statistique multivariée} (GAMMA)
    \item \textbf{interprétation métier} (BA)
  \end{itemize}
\end{block}

\begin{alertblock}{Différence avec le ML supervisé}
  Pas de variable cible $y$ : on cherche à \alert{décrire} et
  \alert{comprendre} la structure, pas à \alert{prédire}.
\end{alertblock}
\end{columns}
\end{frame}

% ────────────────────────────────────────────────────────────
\subsection{1.2 Tableau de données}
% ────────────────────────────────────────────────────────────
\begin{frame}{1.2 — Le tableau de données}
\begin{definition}[Tableau de données]
On note $\mathbf{X} \in \R^{n \times p}$ le tableau dont la ligne $i$ (individu) est
$\mathbf{x}_i = (x_{i1}, \ldots, x_{ip})^\top \in \R^p$, et la colonne $j$
(variable) est $\mathbf{x}^j = (x_{1j}, \ldots, x_{nj})^\top \in \R^n$.
\end{definition}

\begin{columns}[T]
\column{0.50\textwidth}
\textbf{Deux lectures}
\begin{itemize}
  \item \textbf{Ligne} : nuage des $n$ individus dans $\R^p$
  \item \textbf{Colonne} : nuage des $p$ variables dans $\R^n$
\end{itemize}
\vspace{0.2cm}
\textbf{Types de variables}
\begin{itemize}\small
  \item \emph{Quantitatives} (continues, discrètes) : ACP
  \item \emph{Qualitatives nominales} : ACM
  \item \emph{Qualitatives ordinales} : recodage prudent
  \item \emph{Effectifs / contingence} : AFC
\end{itemize}
\column{0.46\textwidth}
\centering
\footnotesize
\renewcommand{\arraystretch}{1.1}
\begin{tabular}{c|ccccc}
  & $V_1$ & $V_2$ & $\cdots$ & $V_p$ \\
\hline
$i_1$    & $x_{11}$ & $x_{12}$ & $\cdots$ & $x_{1p}$ \\
$i_2$    & $x_{21}$ & $x_{22}$ & $\cdots$ & $x_{2p}$ \\
$\vdots$ & $\vdots$ & $\vdots$ & $\ddots$ & $\vdots$ \\
$i_n$    & $x_{n1}$ & $x_{n2}$ & $\cdots$ & $x_{np}$ \\
\end{tabular}

\vspace{0.3cm}
\begin{block}{\small Exemple}
\small Tableau \texttt{decathlon} : $n=41$ athlètes,
$p=10$ disciplines $+$ classement.
\end{block}
\end{columns}
\end{frame}

% ────────────────────────────────────────────────────────────
\subsection{1.3 Statistiques descriptives univariées}
% ────────────────────────────────────────────────────────────
\begin{frame}{1.3 — Caractéristiques d'une variable}
\begin{columns}[T]
\column{0.50\textwidth}
\textbf{Tendance centrale}
\begin{align*}
\bar{x}^j &= \frac{1}{n}\sum_{i=1}^n x_{ij}
  \quad\text{(moyenne arithmétique)}\\
\tilde{x}^j &= \text{médiane}(\mathbf{x}^j)
\end{align*}

\textbf{Dispersion}
\begin{align*}
s_j^2 &= \frac{1}{n}\sum_{i=1}^n (x_{ij} - \bar{x}^j)^2
  \quad\text{(variance empirique)}\\
s_j   &= \sqrt{s_j^2} \quad\text{(écart-type)}
\end{align*}

\begin{rema}
On utilise ici la variance \alert{biaisée} (divisée par $n$),
adaptée à l'analyse géométrique.
\end{rema}
\column{0.46\textwidth}
\textbf{Quantiles et étendue}
\begin{itemize}\small
  \item $Q_1, Q_2 = \tilde{x}, Q_3$ : quartiles
  \item $IQR = Q_3 - Q_1$ : écart interquartile
  \item Étendue : $\max - \min$
\end{itemize}

\textbf{Forme}
\begin{itemize}\small
  \item Asymétrie (skewness) : $\gamma_1 = \E[(X-\mu)^3]/\sigma^3$
  \item Aplatissement (kurtosis) : $\gamma_2 = \E[(X-\mu)^4]/\sigma^4 - 3$
\end{itemize}

\begin{alertblock}{Présence d'outliers}
  Préférer $\tilde{x}$ et $IQR$ à $\bar{x}$ et $s$ pour résumer une
  variable contaminée.
\end{alertblock}
\end{columns}
\end{frame}

% ────────────────────────────────────────────────────────────
\subsection{1.4 Statistiques descriptives bivariées}
% ────────────────────────────────────────────────────────────
\begin{frame}{1.4 — Liaison entre deux variables}
\begin{definition}[Covariance et corrélation empiriques]
Pour deux variables $\mathbf{x}^j, \mathbf{x}^k \in \R^n$ :
\[
\Cov(\mathbf{x}^j, \mathbf{x}^k)
= \frac{1}{n}\sum_{i=1}^n (x_{ij}-\bar{x}^j)(x_{ik}-\bar{x}^k),
\qquad
\Cor(\mathbf{x}^j, \mathbf{x}^k) = \frac{\Cov(\mathbf{x}^j, \mathbf{x}^k)}{s_j \, s_k}.
\]
\end{definition}

\begin{prop}[Propriétés de la corrélation]
\begin{enumerate}
  \item $\Cor(\mathbf{x}^j, \mathbf{x}^k) \in [-1, 1]$ \hfill (Cauchy-Schwarz)
  \item $|\Cor| = 1 \iff \mathbf{x}^k$ est fonction \emph{affine} de $\mathbf{x}^j$
  \item $\Cor$ est invariante par changement d'unités linéaire (échelle, origine)
  \item $\Cor = 0 \not\Rightarrow$ indépendance (mesure uniquement la liaison \emph{linéaire})
\end{enumerate}
\end{prop}

\vspace{0.1cm}
\begin{block}{Forme matricielle}
$\mathbf{V} = \frac{1}{n} \mathbf{X}_c^\top \mathbf{X}_c$ : matrice de variance-covariance ($p\times p$).\\
$\mathbf{R} = \mathbf{D}_{1/s}\, \mathbf{V}\, \mathbf{D}_{1/s}$ : matrice de corrélation, où $\mathbf{D}_{1/s} = \text{diag}(1/s_j)$.
\end{block}
\end{frame}

% ────────────────────────────────────────────────────────────
\subsection{1.5 Centrage et réduction}
% ────────────────────────────────────────────────────────────
\begin{frame}{1.5 — Centrage et réduction}
\begin{columns}[T]
\column{0.52\textwidth}
\textbf{Centrage}
\[
\mathbf{X}_c = \mathbf{X} - \1_n \bar{\mathbf{x}}^\top, \quad
\bar{\mathbf{x}} = \frac{1}{n}\mathbf{X}^\top \1_n
\]
$\Leftrightarrow$ déplacer l'origine au centre de gravité $\mathbf{g}$.

\vspace{0.3cm}
\textbf{Réduction (standardisation)}
\[
\mathbf{Z} = \mathbf{X}_c \, \mathbf{D}_{1/s}, \quad z_{ij} = \frac{x_{ij} - \bar{x}^j}{s_j}
\]
$\Leftrightarrow$ rendre les variables \alert{adimensionnelles}.

\begin{theorem}[Effet sur la variance-covariance]
\[
\mathbf{V}(\mathbf{X}_c) = \mathbf{V}(\mathbf{X}), \qquad
\mathbf{V}(\mathbf{Z}) = \mathbf{R}(\mathbf{X}).
\]
\end{theorem}
\column{0.44\textwidth}
\begin{block}{Pourquoi réduire ?}
\begin{itemize}\small
  \item Variables d'\alert{unités différentes} (kg, EUR, années)
  \item Variances très différentes
  \item Sans réduction, l'ACP est dominée par la variable de plus grande variance
\end{itemize}
\end{block}

\begin{alertblock}{Quand ne pas réduire ?}
\small Lorsque les variables sont
\textbf{de même nature} et de variances comparables
(ex : températures aux mois de l'année).
\end{alertblock}
\end{columns}
\end{frame}

% ────────────────────────────────────────────────────────────
\subsection{1.6 Distance et métrique}
% ────────────────────────────────────────────────────────────
\begin{frame}{1.6 — Distance entre individus}
\begin{definition}[Métrique]
Une \textbf{métrique} sur $\R^p$ est une matrice $\mathbf{M}$ \alert{symétrique
définie positive}. La distance associée entre $\mathbf{x}_i$ et $\mathbf{x}_{i'}$ est :
\[
d_{\mathbf{M}}^2(\mathbf{x}_i, \mathbf{x}_{i'}) = (\mathbf{x}_i - \mathbf{x}_{i'})^\top \mathbf{M} (\mathbf{x}_i - \mathbf{x}_{i'}).
\]
\end{definition}

\begin{columns}[T]
\column{0.49\textwidth}
\textbf{Cas particuliers}
\begin{itemize}\small
  \item $\mathbf{M} = \mathbf{I}_p$ : distance \alert{euclidienne classique}\\
        $d^2 = \sum_j (x_{ij}-x_{i'j})^2$
  \item $\mathbf{M} = \mathbf{D}_{1/s^2}$ : distance \alert{euclidienne réduite}\\
        $d^2 = \sum_j \frac{(x_{ij}-x_{i'j})^2}{s_j^2}$
  \item $\mathbf{M} = \mathbf{V}^{-1}$ : \alert{Mahalanobis}
  \item $\mathbf{M} = \mathbf{D}_{1/x_{\bullet j}}$ : $\chi^2$ (AFC)
\end{itemize}
\column{0.49\textwidth}
\begin{prop}[Choix de la métrique]
Le choix de $\mathbf{M}$ \emph{définit} l'analyse :
\begin{itemize}\small
  \item ACP normée $\equiv$ ACP centrée avec $\mathbf{M} = \mathbf{D}_{1/s^2}$
  \item AFC : tableau de profils, $\mathbf{M}_{\chi^2}$
  \item AFD : $\mathbf{M} = \mathbf{V}_{\text{intra}}^{-1}$
\end{itemize}
\end{prop}

\begin{block}{Pondération des individus}
\small Chaque individu a un poids $p_i \geq 0$ tel que
$\sum_i p_i = 1$. Cas usuel : $p_i = 1/n$.
\end{block}
\end{columns}
\end{frame}

% ────────────────────────────────────────────────────────────
\subsection{1.7 Notion d'inertie}
% ────────────────────────────────────────────────────────────
\begin{frame}{1.7 — Inertie : la mesure de l'analyse des données}
\begin{definition}[Inertie totale d'un nuage]
Soit $\mathcal{N} = \{(\mathbf{x}_i, p_i)\}_{i=1}^n$ un nuage pondéré de centre de gravité
$\mathbf{g} = \sum_i p_i \mathbf{x}_i$. L'\textbf{inertie totale} de $\mathcal{N}$ par rapport à
$\mathbf{g}$, dans la métrique $\mathbf{M}$, est :
\[
I_{\mathcal{N}} = \sum_{i=1}^n p_i \, d_{\mathbf{M}}^2(\mathbf{x}_i, \mathbf{g})
= \sum_{i=1}^n p_i \, (\mathbf{x}_i - \mathbf{g})^\top \mathbf{M} (\mathbf{x}_i - \mathbf{g}).
\]
\end{definition}

\begin{prop}[Lien avec la trace]
Soit $\mathbf{V} = \sum_i p_i (\mathbf{x}_i - \mathbf{g})(\mathbf{x}_i - \mathbf{g})^\top$
la matrice de variance-covariance pondérée. Alors :
\[
\boxed{\,I_{\mathcal{N}} = \Tr(\mathbf{V} \mathbf{M}).\,}
\]
\end{prop}

\begin{rema}
Pour $\mathbf{M} = \mathbf{I}_p$ et $p_i = 1/n$ : $I_{\mathcal{N}} = \Tr(\mathbf{V}) = \sum_j s_j^2$,
soit la \textbf{somme des variances} des variables. En ACP normée :
$I_{\mathcal{N}} = \Tr(\mathbf{R}) = p$.
\end{rema}
\end{frame}

% ────────────────────────────────────────────────────────────
\subsection{1.8 Théorème d'Huygens}
% ────────────────────────────────────────────────────────────
\begin{frame}{1.8 — Théorème de Huygens (décomposition de l'inertie)}
\begin{theorem}[Huygens]
Soit $\mathcal{N}$ un nuage partitionné en $K$ classes $\mathcal{C}_1, \ldots, \mathcal{C}_K$
de poids respectifs $P_k = \sum_{i \in \mathcal{C}_k} p_i$ et de barycentres
$\mathbf{g}_k = \frac{1}{P_k}\sum_{i \in \mathcal{C}_k} p_i \mathbf{x}_i$. Alors :
\[
\boxed{\;
\underbrace{\sum_{i} p_i \, d^2(\mathbf{x}_i, \mathbf{g})}_{I_{\text{totale}}}
=
\underbrace{\sum_{k=1}^K \sum_{i \in \mathcal{C}_k} p_i \, d^2(\mathbf{x}_i, \mathbf{g}_k)}_{I_{\text{intra}}}
\;+\;
\underbrace{\sum_{k=1}^K P_k \, d^2(\mathbf{g}_k, \mathbf{g})}_{I_{\text{inter}}}.
\;}
\]
\end{theorem}

\begin{columns}[T]
\column{0.55\textwidth}
\textbf{Conséquences fondamentales}
\begin{itemize}\small
  \item \alert{Clustering} : minimiser $I_{\text{intra}}$ $\Leftrightarrow$
        maximiser $I_{\text{inter}}$ (k-means)
  \item \alert{Réduction de dimension} : maximiser
        l'inertie projetée $\Leftrightarrow$ minimiser
        l'inertie résiduelle (ACP)
  \item \alert{AFD} : critère de Fisher
        $\max \, I_{\text{inter}} / I_{\text{intra}}$
\end{itemize}
\column{0.42\textwidth}
\begin{tikzpicture}[scale=0.9]
\fill[deepblue!10] (0.5,0.5) circle(0.5cm);
\fill[deeporange!10] (2.5,0.7) circle(0.5cm);
\fill[darkgreen!10] (1.7,2.0) circle(0.45cm);
\foreach \x/\y in {0.2/0.5, 0.4/0.7, 0.7/0.4, 0.5/0.3}
  \fill[deepblue] (\x,\y) circle(1.3pt);
\foreach \x/\y in {2.3/0.6, 2.6/0.8, 2.7/0.5, 2.4/0.9}
  \fill[deeporange] (\x,\y) circle(1.3pt);
\foreach \x/\y in {1.6/2.0, 1.8/2.1, 1.7/1.8}
  \fill[darkgreen] (\x,\y) circle(1.3pt);
\fill[deepblue] (0.45,0.5) circle(2.5pt) node[left, font=\tiny] {$\mathbf{g}_1$};
\fill[deeporange] (2.5,0.7) circle(2.5pt) node[right, font=\tiny] {$\mathbf{g}_2$};
\fill[darkgreen] (1.7,2.0) circle(2.5pt) node[above, font=\tiny] {$\mathbf{g}_3$};
\fill[black] (1.55,1.06) circle(2.5pt) node[right, font=\tiny] {$\mathbf{g}$};
\draw[dashed, gray] (0.45,0.5) -- (1.55,1.06);
\draw[dashed, gray] (2.5,0.7)  -- (1.55,1.06);
\draw[dashed, gray] (1.7,2.0)  -- (1.55,1.06);
\end{tikzpicture}
\end{columns}
\end{frame}

% ────────────────────────────────────────────────────────────
\subsection{1.9 Le tableau comme deux nuages}
% ────────────────────────────────────────────────────────────
\begin{frame}{1.9 — Dualité individus / variables}
\begin{columns}[T]
\column{0.49\textwidth}
\textbf{Nuage des individus $\mathcal{N}_I$ dans $\R^p$}
\begin{itemize}\small
  \item $n$ points $\mathbf{x}_i$
  \item Métrique $\mathbf{M}$ sur $\R^p$
  \item Pondération $\mathbf{D}_p = \text{diag}(p_i)$
  \item Centre $\mathbf{g}$, inertie $\Tr(\mathbf{V}\mathbf{M})$
\end{itemize}

\vspace{0.2cm}
\textbf{Nuage des variables $\mathcal{N}_V$ dans $\R^n$}
\begin{itemize}\small
  \item $p$ points $\mathbf{x}^j$
  \item Métrique $\mathbf{D}_p$ sur $\R^n$
  \item Pondération $\mathbf{M}$ sur les variables
  \item Produit scalaire $\langle \mathbf{x}^j, \mathbf{x}^k \rangle = \Cov(\mathbf{x}^j, \mathbf{x}^k)$
\end{itemize}

\column{0.49\textwidth}
\begin{theorem}[Dualité]
Les deux nuages $\mathcal{N}_I$ et $\mathcal{N}_V$ partagent les
\alert{mêmes valeurs propres non nulles}. La même décomposition
matricielle (SVD de $\mathbf{X}$) fournit simultanément :
\begin{itemize}\small
  \item les axes principaux de $\mathcal{N}_I$
  \item les axes principaux de $\mathcal{N}_V$
\end{itemize}
\end{theorem}

\begin{block}{Conséquence pratique}
\small Les \textbf{individus} et les \textbf{variables} se représentent
sur les \emph{mêmes} plans factoriels (biplot) — l'ACP est
\alert{auto-duale}.
\end{block}
\end{columns}
\end{frame}

% ────────────────────────────────────────────────────────────
\subsection{1.10 Aperçu R}
% ────────────────────────────────────────────────────────────
\begin{frame}[fragile]{1.10 — Aperçu : statistiques descriptives en R}
\begin{lstlisting}
library(FactoMineR)    # methodes factorielles
library(factoextra)    # visualisation
library(corrplot)      # matrice de correlations

data(decathlon)        # 41 athletes x 13 variables
X <- decathlon[, 1:10] # variables actives quantitatives

# Statistiques univariees
summary(X)
sapply(X, sd)

# Centrage / reduction
Z <- scale(X)          # par defaut center=TRUE, scale=TRUE
colMeans(Z); apply(Z, 2, sd)

# Matrice de variance-covariance et de correlation
V <- cov(X)            # divise par n-1 (correction de Bessel)
R <- cor(X)
corrplot(R, method = "color", type = "upper", tl.cex = 0.7)

# Inertie totale (ACP normee)
sum(diag(R))           # = nombre de variables = 10
\end{lstlisting}
\end{frame}

% ────────────────────────────────────────────────────────────
\subsection{1.11 Récapitulatif}
% ────────────────────────────────────────────────────────────
\begin{frame}{1.11 — Récapitulatif du chapitre}
\begin{columns}[T]
\column{0.49\textwidth}
\textbf{Concepts à retenir}
\begin{itemize}\small
  \item Tableau $\mathbf{X} \in \R^{n\times p}$, deux lectures
        (individus / variables)
  \item Matrice de variance-covariance
        $\mathbf{V} = \frac{1}{n}\mathbf{X}_c^\top \mathbf{X}_c$
  \item Matrice de corrélation $\mathbf{R}$
  \item Centrage (décalage), réduction (échelle)
  \item Métrique $\mathbf{M}$ symétrique définie positive
  \item Inertie $I = \Tr(\mathbf{V}\mathbf{M})$
  \item Théorème de Huygens : $I = I_{\text{intra}} + I_{\text{inter}}$
\end{itemize}

\column{0.49\textwidth}
\begin{block}{À venir}
\textbf{Chapitre 2} : algèbre linéaire pour l'analyse des données
\begin{itemize}\small
  \item Théorème spectral
  \item Décomposition en valeurs singulières (SVD)
  \item Théorème d'Eckart-Young
  \item Projections orthogonales
\end{itemize}
\end{block}

\begin{alertblock}{Idée directrice}
Toute l'analyse des données est une
\textbf{décomposition d'inertie} via la diagonalisation
d'une matrice symétrique.
\end{alertblock}
\end{columns}
\end{frame}

% ============================================================
%  Chapitre 2 — Algèbre linéaire pour l'analyse des données
% ============================================================
\section{Algèbre Linéaire pour l'Analyse des Données}

% ────────────────────────────────────────────────────────────
\subsection{2.1 Rappels : valeurs et vecteurs propres}
% ────────────────────────────────────────────────────────────
\begin{frame}{2.1 — Rappels d'algèbre linéaire}
\begin{definition}[Valeur propre, vecteur propre]
Soit $\mathbf{A} \in \R^{p \times p}$. Un scalaire $\lambda \in \R$ est dit
\textbf{valeur propre} de $\mathbf{A}$ s'il existe $\mathbf{u} \in \R^p \setminus \{0\}$
tel que $\mathbf{A}\mathbf{u} = \lambda \mathbf{u}$. Le vecteur $\mathbf{u}$ est
le \textbf{vecteur propre} associé.
\end{definition}

\begin{columns}[T]
\column{0.49\textwidth}
\textbf{Polynôme caractéristique}
\[
P_{\mathbf{A}}(\lambda) = \det(\mathbf{A} - \lambda \mathbf{I}_p)
\]
Les valeurs propres sont les racines de $P_{\mathbf{A}}$ dans $\mathbb{C}$.
\vspace{0.2cm}

\textbf{Spectre} : $\sigma(\mathbf{A}) = \{\lambda_1, \ldots, \lambda_p\}$
(comptées avec multiplicité).

\textbf{Sous-espace propre}
\[
E_\lambda = \ker(\mathbf{A} - \lambda \mathbf{I}_p)
\]

\column{0.49\textwidth}
\begin{proposition}[Trace et déterminant]
\[
\Tr(\mathbf{A}) = \sum_{i=1}^p \lambda_i, \qquad
\det(\mathbf{A}) = \prod_{i=1}^p \lambda_i.
\]
\end{proposition}

\begin{proposition}[Diagonalisation]
$\mathbf{A}$ est \alert{diagonalisable} si et seulement si la somme des
dimensions des sous-espaces propres vaut $p$ :
\[
\sum_{\lambda \in \sigma(\mathbf{A})} \dim(E_\lambda) = p.
\]
On peut alors écrire $\mathbf{A} = \mathbf{P}\mathbf{D}\mathbf{P}^{-1}$
avec $\mathbf{D}$ diagonale.
\end{proposition}
\end{columns}
\end{frame}

% ────────────────────────────────────────────────────────────
\subsection{2.2 Théorème spectral}
% ────────────────────────────────────────────────────────────
\begin{frame}{2.2 — Théorème spectral pour matrices symétriques}
\begin{theorem}[Théorème spectral]
Soit $\mathbf{A} \in \R^{p \times p}$ une matrice \alert{symétrique réelle}
($\mathbf{A}^\top = \mathbf{A}$). Alors :
\begin{enumerate}
  \item Toutes les valeurs propres de $\mathbf{A}$ sont \textbf{réelles}.
  \item Les sous-espaces propres associés à des valeurs propres distinctes
        sont \textbf{orthogonaux}.
  \item Il existe une \textbf{base orthonormée} $\{\mathbf{u}_1, \ldots, \mathbf{u}_p\}$
        de $\R^p$ formée de vecteurs propres de $\mathbf{A}$.
\end{enumerate}
On a la décomposition :
\[
\boxed{\;\mathbf{A} = \mathbf{U}\, \mathbf{\Lambda}\, \mathbf{U}^\top
       = \sum_{k=1}^p \lambda_k \mathbf{u}_k \mathbf{u}_k^\top \;}
\]
avec $\mathbf{U}$ orthogonale ($\mathbf{U}^\top \mathbf{U} = \mathbf{I}_p$) et
$\mathbf{\Lambda} = \text{diag}(\lambda_1, \ldots, \lambda_p)$.
\end{theorem}

\begin{remark}[Conséquence majeure]
\small La matrice de variance-covariance $\mathbf{V}$ et la matrice de corrélation
$\mathbf{R}$ sont symétriques réelles : elles admettent toujours une base
\alert{orthonormée} de vecteurs propres. C'est le \textbf{cœur} de l'ACP.
\end{remark}
\end{frame}

% ────────────────────────────────────────────────────────────
\subsection{2.3 Formes quadratiques}
% ────────────────────────────────────────────────────────────
\begin{frame}{2.3 — Formes quadratiques et matrices SDP}
\begin{definition}[Forme quadratique]
À toute matrice symétrique $\mathbf{A} \in \R^{p\times p}$ on associe la
\textbf{forme quadratique}
\[
q_{\mathbf{A}} : \R^p \to \R, \quad
q_{\mathbf{A}}(\mathbf{x}) = \mathbf{x}^\top \mathbf{A}\, \mathbf{x}.
\]
\end{definition}

\begin{definition}[Matrice positive / définie positive]
\begin{itemize}
  \item $\mathbf{A}$ est \textbf{positive} ($\mathbf{A} \succeq 0$) si $q_{\mathbf{A}}(\mathbf{x}) \geq 0$ pour tout $\mathbf{x}$.
  \item $\mathbf{A}$ est \textbf{définie positive} ($\mathbf{A} \succ 0$) si $q_{\mathbf{A}}(\mathbf{x}) > 0$ pour tout $\mathbf{x} \neq 0$.
\end{itemize}
\end{definition}

\begin{theorem}[Caractérisation par les valeurs propres]
Soit $\mathbf{A}$ symétrique réelle. Alors :
\[
\mathbf{A} \succeq 0 \iff \forall k, \; \lambda_k \geq 0,
\qquad
\mathbf{A} \succ 0 \iff \forall k, \; \lambda_k > 0.
\]
\end{theorem}

\begin{remark}
$\mathbf{V}$ et $\mathbf{R}$ sont \alert{toujours} positives car
$\mathbf{V} = \frac{1}{n}\mathbf{X}_c^\top \mathbf{X}_c$ : pour tout $\mathbf{u}$,
$\mathbf{u}^\top \mathbf{V} \mathbf{u} = \frac{1}{n}\|\mathbf{X}_c \mathbf{u}\|^2 \geq 0$.
\end{remark}
\end{frame}

% ────────────────────────────────────────────────────────────
\subsection{2.4 Quotient de Rayleigh}
% ────────────────────────────────────────────────────────────
\begin{frame}{2.4 — Quotient de Rayleigh}
\begin{definition}[Quotient de Rayleigh]
Soit $\mathbf{A}$ symétrique réelle et $\mathbf{u} \in \R^p \setminus \{0\}$. Le
\textbf{quotient de Rayleigh} est :
\[
R_{\mathbf{A}}(\mathbf{u}) = \frac{\mathbf{u}^\top \mathbf{A}\, \mathbf{u}}{\mathbf{u}^\top \mathbf{u}}.
\]
\end{definition}

\begin{theorem}[Courant-Fischer-Weyl]
Soient $\lambda_1 \geq \lambda_2 \geq \cdots \geq \lambda_p$ les valeurs propres
de $\mathbf{A}$ (symétrique). Alors :
\[
\boxed{\;\lambda_1 = \max_{\mathbf{u} \neq 0} R_{\mathbf{A}}(\mathbf{u}),
\qquad
\lambda_p = \min_{\mathbf{u} \neq 0} R_{\mathbf{A}}(\mathbf{u}).\;}
\]
Plus généralement, pour le $k$-ième axe :
\[
\lambda_k = \max_{\substack{\mathbf{u} \perp \mathbf{u}_1, \ldots, \mathbf{u}_{k-1} \\ \mathbf{u} \neq 0}}
R_{\mathbf{A}}(\mathbf{u}).
\]
\end{theorem}

\begin{remark}[Importance pour l'ACP]
\small Maximiser la variance projetée $\mathbf{u}^\top \mathbf{V} \mathbf{u}$ sous
$\|\mathbf{u}\|=1$ revient à trouver le \alert{vecteur propre dominant} de $\mathbf{V}$.
C'est l'optimisation centrale de l'ACP.
\end{remark}
\end{frame}

% ────────────────────────────────────────────────────────────
\subsection{2.5 Décomposition en valeurs singulières}
% ────────────────────────────────────────────────────────────
\begin{frame}{2.5 — Décomposition en valeurs singulières (SVD)}
\begin{theorem}[SVD]
Toute matrice $\mathbf{X} \in \R^{n \times p}$ admet une décomposition
\[
\boxed{\;\mathbf{X} = \mathbf{U}\, \mathbf{\Sigma}\, \mathbf{V}^\top\;}
\]
où :
\begin{itemize}
  \item $\mathbf{U} \in \R^{n \times n}$ orthogonale ($\mathbf{U}^\top \mathbf{U} = \mathbf{I}_n$)
  \item $\mathbf{V} \in \R^{p \times p}$ orthogonale ($\mathbf{V}^\top \mathbf{V} = \mathbf{I}_p$)
  \item $\mathbf{\Sigma} \in \R^{n \times p}$ rectangulaire diagonale, de coefficients
        diagonaux $\sigma_1 \geq \sigma_2 \geq \cdots \geq \sigma_r \geq 0$,
        où $r = \rang(\mathbf{X})$.
\end{itemize}
Les $\sigma_k$ sont les \textbf{valeurs singulières}, les colonnes de
$\mathbf{U}$ et $\mathbf{V}$ sont les \textbf{vecteurs singuliers à gauche} et
\textbf{à droite}.
\end{theorem}

\begin{proposition}[Forme dyadique / SVD réduite]
\[
\mathbf{X} = \sum_{k=1}^r \sigma_k\, \mathbf{u}_k \mathbf{v}_k^\top.
\]
Chaque terme est une matrice de rang 1. La SVD décompose $\mathbf{X}$ en
\alert{somme} de matrices de rang 1, par ordre décroissant d'importance.
\end{proposition}
\end{frame}

% ────────────────────────────────────────────────────────────
\subsection{2.6 Lien spectral / SVD}
% ────────────────────────────────────────────────────────────
\begin{frame}{2.6 — Lien entre SVD et décomposition spectrale}
\begin{theorem}[Spectres de $\mathbf{X}^\top \mathbf{X}$ et $\mathbf{X} \mathbf{X}^\top$]
Soit $\mathbf{X} = \mathbf{U}\mathbf{\Sigma}\mathbf{V}^\top$ la SVD de $\mathbf{X}$. Alors :
\[
\mathbf{X}^\top \mathbf{X} = \mathbf{V}\, \mathbf{\Sigma}^\top \mathbf{\Sigma}\, \mathbf{V}^\top,
\qquad
\mathbf{X} \mathbf{X}^\top = \mathbf{U}\, \mathbf{\Sigma} \mathbf{\Sigma}^\top\, \mathbf{U}^\top.
\]
Les valeurs propres communes (non nulles) sont $\lambda_k = \sigma_k^2$.
\end{theorem}

\begin{columns}[T]
\column{0.49\textwidth}
\textbf{Démonstration (idée).}
\[
\mathbf{X}^\top \mathbf{X}
= (\mathbf{U}\mathbf{\Sigma}\mathbf{V}^\top)^\top (\mathbf{U}\mathbf{\Sigma}\mathbf{V}^\top)
= \mathbf{V}\mathbf{\Sigma}^\top \mathbf{U}^\top \mathbf{U} \mathbf{\Sigma}\mathbf{V}^\top
\]
$= \mathbf{V}\, (\mathbf{\Sigma}^\top \mathbf{\Sigma})\, \mathbf{V}^\top$
car $\mathbf{U}^\top \mathbf{U} = \mathbf{I}_n$. La matrice $\mathbf{\Sigma}^\top \mathbf{\Sigma}$
est diagonale de coefficients $\sigma_k^2$.

\column{0.49\textwidth}
\begin{block}{Conséquence pour l'ACP}
\small La SVD de $\mathbf{X}_c$ donne directement :
\begin{itemize}\small
  \item Les axes principaux : colonnes de $\mathbf{V}$
  \item Les composantes principales : $\mathbf{C} = \mathbf{X}_c \mathbf{V}$
  \item Les inerties expliquées : $\lambda_k = \sigma_k^2$
\end{itemize}
\end{block}

\begin{alertblock}{En pratique}
\small On utilise la \textbf{SVD numérique} pour calculer l'ACP, plus stable
que la diagonalisation directe de $\mathbf{V}$.
\end{alertblock}
\end{columns}
\end{frame}

% ────────────────────────────────────────────────────────────
\subsection{2.7 Théorème d'Eckart-Young}
% ────────────────────────────────────────────────────────────
\begin{frame}{2.7 — Théorème d'Eckart-Young-Mirsky}
\begin{theorem}[Eckart-Young, 1936]
Soit $\mathbf{X} = \sum_{k=1}^r \sigma_k \mathbf{u}_k \mathbf{v}_k^\top$ la SVD de $\mathbf{X}$.
Pour tout entier $q \leq r$, la matrice
\[
\mathbf{X}_q = \sum_{k=1}^q \sigma_k \, \mathbf{u}_k \mathbf{v}_k^\top
\]
est l'unique matrice de rang $\leq q$ minimisant l'erreur de Frobenius :
\[
\mathbf{X}_q = \argmin_{\rang(\mathbf{Y}) \leq q} \|\mathbf{X} - \mathbf{Y}\|_F^2,
\quad
\text{avec}
\quad
\|\mathbf{X} - \mathbf{X}_q\|_F^2 = \sum_{k=q+1}^{r} \sigma_k^2.
\]
\end{theorem}

\begin{block}{Norme de Frobenius}
\small $\displaystyle \|\mathbf{A}\|_F^2 = \sum_{i,j} a_{ij}^2 = \Tr(\mathbf{A}^\top \mathbf{A}) = \sum_k \sigma_k^2$.
\end{block}

\begin{remark}[Interprétation pour l'analyse des données]
\small Eckart-Young fonde \alert{toute la réduction de dimension} : la
\textbf{meilleure approximation de rang $q$} d'un nuage de points est la
projection sur le sous-espace engendré par les $q$ premiers vecteurs propres
de $\mathbf{X}_c^\top \mathbf{X}_c$ — c'est l'ACP.
\end{remark}
\end{frame}

% ────────────────────────────────────────────────────────────
\subsection{2.8 Projections orthogonales}
% ────────────────────────────────────────────────────────────
\begin{frame}{2.8 — Projections orthogonales}
\begin{definition}[Projecteur orthogonal]
Soit $F$ un sous-espace vectoriel de $\R^p$. Le \textbf{projecteur orthogonal}
sur $F$ est l'unique application linéaire $\mathbf{P}_F : \R^p \to F$ telle que
$\mathbf{x} - \mathbf{P}_F(\mathbf{x}) \perp F$ pour tout $\mathbf{x}$.
Il vérifie : $\mathbf{P}_F^2 = \mathbf{P}_F$ et $\mathbf{P}_F^\top = \mathbf{P}_F$.
\end{definition}

\begin{proposition}[Forme matricielle]
Si $\{\mathbf{u}_1, \ldots, \mathbf{u}_q\}$ est une base orthonormée de $F$ et
$\mathbf{U} = [\mathbf{u}_1 \,|\, \cdots \,|\, \mathbf{u}_q]$, alors :
\[
\mathbf{P}_F = \mathbf{U} \mathbf{U}^\top.
\]
La projection de $\mathbf{x}$ est $\mathbf{P}_F(\mathbf{x}) = \mathbf{U}(\mathbf{U}^\top \mathbf{x}) = \sum_{k=1}^q \langle \mathbf{x}, \mathbf{u}_k\rangle \mathbf{u}_k$.
\end{proposition}

\begin{theorem}[Théorème de Pythagore]
Pour tout $\mathbf{x} \in \R^p$ :
\[
\|\mathbf{x}\|^2 = \|\mathbf{P}_F(\mathbf{x})\|^2 + \|\mathbf{x} - \mathbf{P}_F(\mathbf{x})\|^2.
\]
\end{theorem}

\begin{remark}
\small En analyse des données : projeter sur l'espace propre des $q$ premiers axes
de $\mathbf{V}$ \alert{maximise} la norme projetée donc l'inertie projetée.
\end{remark}
\end{frame}

% ────────────────────────────────────────────────────────────
\subsection{2.9 Inertie projetée}
% ────────────────────────────────────────────────────────────
\begin{frame}{2.9 — Inertie projetée et meilleur sous-espace}
\begin{theorem}[Décomposition de l'inertie]
Soit $\mathcal{N}$ un nuage centré de matrice de variance-covariance $\mathbf{V}$,
et $F_q$ le sous-espace engendré par les $q$ premiers vecteurs propres de $\mathbf{V}$.
Alors l'inertie totale se décompose comme :
\[
\boxed{\;
I_{\text{totale}}
= \underbrace{\sum_{k=1}^q \lambda_k}_{\text{inertie projetée sur } F_q}
+ \underbrace{\sum_{k=q+1}^p \lambda_k}_{\text{inertie résiduelle}}.
\;}
\]
$F_q$ est le sous-espace de dimension $q$ \alert{maximisant} l'inertie projetée.
\end{theorem}

\begin{proposition}[Pourcentage d'inertie expliquée]
La part d'inertie expliquée par les $q$ premiers axes est :
\[
\tau_q = \frac{\lambda_1 + \cdots + \lambda_q}{\lambda_1 + \cdots + \lambda_p} \in [0, 1].
\]
On choisit $q$ pour que $\tau_q$ atteigne un seuil (ex : 80\%).
\end{proposition}
\end{frame}

% ────────────────────────────────────────────────────────────
\subsection{2.10 Métriques généralisées}
% ────────────────────────────────────────────────────────────
\begin{frame}{2.10 — Décomposition généralisée (avec métrique $\mathbf{M}$)}
\begin{theorem}[SVD généralisée]
Soient $\mathbf{M}$ une métrique sur $\R^p$ et $\mathbf{D}_p$ une métrique sur $\R^n$
(symétriques définies positives). Il existe une décomposition de $\mathbf{X}$
\[
\mathbf{X} = \mathbf{U}\, \mathbf{\Sigma}\, \mathbf{V}^\top
\]
avec $\mathbf{U}^\top \mathbf{D}_p \mathbf{U} = \mathbf{I}$ et $\mathbf{V}^\top \mathbf{M} \mathbf{V} = \mathbf{I}$.
Les vecteurs propres $\mathbf{v}_k$ diagonalisent $\mathbf{X}^\top \mathbf{D}_p \mathbf{X} \mathbf{M}$.
\end{theorem}

\begin{columns}[T]
\column{0.49\textwidth}
\textbf{Cadre unificateur}\\
\small Toutes les méthodes factorielles s'écrivent comme une SVD
généralisée d'un triplet $(\mathbf{X}, \mathbf{M}, \mathbf{D}_p)$ :

\begin{itemize}\small
  \item ACP : $\mathbf{M} = \mathbf{I}_p$ ou $\mathbf{D}_{1/s^2}$, $\mathbf{D}_p = \frac{1}{n}\mathbf{I}_n$
  \item AFC : $\mathbf{M} = \mathbf{D}_{1/x_{\bullet j}}$, $\mathbf{D}_p$ liée aux profils
  \item ACM : AFC du tableau disjonctif
  \item AFD : $\mathbf{M} = \mathbf{V}_{\text{intra}}^{-1}$
\end{itemize}
\column{0.49\textwidth}
\begin{block}{Vue duale}
\small Le \textbf{schéma de dualité} (Cailliez-Pagès) :
\[
\mathbf{X}^\top \mathbf{D}_p \mathbf{X} \mathbf{M}
\quad \stackrel{\text{spectres}}{\equiv} \quad
\mathbf{X} \mathbf{M} \mathbf{X}^\top \mathbf{D}_p
\]
Mêmes valeurs propres non nulles : c'est la dualité individus/variables.
\end{block}

\begin{alertblock}{Conséquence}
\small Une seule machine algébrique --- la SVD généralisée --- couvre
l'\emph{ensemble} des méthodes factorielles classiques.
\end{alertblock}
\end{columns}
\end{frame}

% ────────────────────────────────────────────────────────────
\subsection{2.11 Aperçu R}
% ────────────────────────────────────────────────────────────
\begin{frame}[fragile]{2.11 — Algèbre linéaire en R}
\begin{lstlisting}
# Matrice symetrique definie positive
set.seed(42)
A <- matrix(rnorm(16), 4, 4); A <- t(A) %*% A   # symetrique, SDP

# Decomposition spectrale
res <- eigen(A)
res$values    # valeurs propres (decroissantes)
res$vectors   # vecteurs propres en colonnes (orthonormes)

# Verification : A = U D U^T
U <- res$vectors; D <- diag(res$values)
all.equal(A, U %*% D %*% t(U))   # TRUE

# SVD d'une matrice generale
X <- matrix(rnorm(20), 5, 4)
svd_res <- svd(X)
svd_res$d     # valeurs singulieres
svd_res$u; svd_res$v

# Verification : X = U Sigma V^T
all.equal(X, svd_res$u %*% diag(svd_res$d) %*% t(svd_res$v))  # TRUE

# Approximation de rang 2 (Eckart-Young)
X2 <- svd_res$u[, 1:2] %*% diag(svd_res$d[1:2]) %*% t(svd_res$v[, 1:2])
norm(X - X2, type = "F")^2          # = somme des sigma^2 restants
sum(svd_res$d[3:4]^2)
\end{lstlisting}
\end{frame}

% ────────────────────────────────────────────────────────────
\subsection{2.12 Récapitulatif}
% ────────────────────────────────────────────────────────────
\begin{frame}{2.12 — Récapitulatif du chapitre}
\begin{columns}[T]
\column{0.49\textwidth}
\textbf{Théorèmes à retenir}
\begin{itemize}\small
  \item \textbf{Spectral} : $\mathbf{A}$ symétrique réelle
        $\Rightarrow$ base orthonormée de vecteurs propres
        $\mathbf{A} = \mathbf{U}\mathbf{\Lambda}\mathbf{U}^\top$
  \item \textbf{Courant-Fischer-Weyl} : $\lambda_1 = \max R_{\mathbf{A}}(\mathbf{u})$
  \item \textbf{SVD} : $\mathbf{X} = \mathbf{U}\mathbf{\Sigma}\mathbf{V}^\top$ pour toute matrice
  \item \textbf{Eckart-Young} : meilleure approximation de rang $q$
  \item \textbf{Pythagore} : $\|\mathbf{x}\|^2 = \|\mathbf{P}_F\mathbf{x}\|^2 + \|\mathbf{x} - \mathbf{P}_F\mathbf{x}\|^2$
\end{itemize}

\column{0.49\textwidth}
\begin{block}{Pont vers l'ACP}
Toutes les briques sont en place :
\begin{enumerate}\small
  \item Variance-covariance $\mathbf{V}$ symétrique SDP
  \item Théorème spectral $\Rightarrow$ axes propres orthogonaux
  \item Quotient de Rayleigh $\Rightarrow$ inertie projetée maximale
  \item Eckart-Young $\Rightarrow$ ACP est optimale
\end{enumerate}
\end{block}

\begin{alertblock}{Idée maîtresse}
L'ACP, l'AFC, l'ACM, l'AFD sont \alert{une seule et même opération} :
diagonaliser une matrice symétrique pour décomposer une inertie.
\end{alertblock}
\end{columns}
\end{frame}

%% Chapitre — Analyse en Composantes Principales (ACP)
%% À compléter
\section{Analyse en Composantes Principales (ACP)}

% ============================================================
%  Chapitre 4 — Analyse Factorielle des Correspondances (AFC)
% ============================================================
\section{Analyse Factorielle des Correspondances (AFC)}

% ────────────────────────────────────────────────────────────
\subsection{4.1 Position du problème}
% ────────────────────────────────────────────────────────────
\begin{frame}{4.1 — Pourquoi l'AFC ?}
\begin{columns}[T]
\column{0.55\textwidth}
\textbf{Données} : table de contingence $\mathbf{N} = (n_{ij}) \in \N^{I \times J}$
croisant deux variables \alert{qualitatives}.\\[0.1cm]

\textbf{Objectifs}
\begin{enumerate}
  \item Décrire les associations entre modalités lignes et colonnes
  \item Tester / quantifier l'\textbf{écart à l'indépendance}
  \item Représenter \textbf{simultanément} lignes et colonnes
        sur un plan factoriel
\end{enumerate}

\vspace{0.2cm}
\textbf{Origine} : Benzécri (1973), école française.
\textbf{AFC} = ACP \emph{généralisée} avec :
\begin{itemize}\small
  \item Tableau de profils
  \item Métrique du $\chi^2$
\end{itemize}
\column{0.42\textwidth}
\centering
\footnotesize
\textbf{CSP $\times$ Diplôme} (INSEE)
\renewcommand{\arraystretch}{1.05}
\begin{tabular}{l|cccc|c}
\toprule
 & Bac & Lic. & Mast. & Doct. & total \\
\midrule
Cadre   & 12 & 65 & 180 & 95 & 352 \\
ProfInt & 88 & 110 & 50  & 5  & 253 \\
Employé & 230& 45 & 12  & 1  & 288 \\
Ouvrier & 195& 8  & 2   & 0  & 205 \\
\midrule
total   & 525& 228& 244 & 101& 1098\\
\bottomrule
\end{tabular}

\vspace{0.4cm}
\begin{block}{\small Question}
\small Quels diplômes sont \emph{associés} à quelles CSP ?
\end{block}
\end{columns}
\end{frame}

% ────────────────────────────────────────────────────────────
\subsection{4.2 Profils lignes et colonnes}
% ────────────────────────────────────────────────────────────
\begin{frame}{4.2 — Profils lignes et profils colonnes}
\begin{definition}[Tableau de fréquences relatives]
\[
f_{ij} = \frac{n_{ij}}{n}, \qquad
f_{i \bullet} = \sum_j f_{ij}, \qquad
f_{\bullet j} = \sum_i f_{ij}
\]
où $n = \sum_{i,j} n_{ij}$ est l'effectif total.
\end{definition}

\begin{definition}[Profil ligne, profil colonne]
\begin{itemize}
  \item \textbf{Profil de la ligne $i$} :
        $\mathbf{L}_i = \left(\dfrac{f_{i1}}{f_{i\bullet}}, \ldots, \dfrac{f_{iJ}}{f_{i\bullet}}\right) \in \R^J$
  \item \textbf{Profil de la colonne $j$} :
        $\mathbf{C}_j = \left(\dfrac{f_{1j}}{f_{\bullet j}}, \ldots, \dfrac{f_{Ij}}{f_{\bullet j}}\right) \in \R^I$
\end{itemize}
Chaque profil est une \alert{distribution conditionnelle} (somme $= 1$).
\end{definition}

\begin{rema}[Profil moyen]
\small Le \textbf{profil moyen ligne} est $(f_{\bullet 1}, \ldots, f_{\bullet J})$ (la marge colonne).
Sous l'\textbf{indépendance} $\chi^2$, tous les profils lignes seraient égaux à ce profil moyen.
\end{rema}
\end{frame}

% ────────────────────────────────────────────────────────────
\subsection{4.3 Distance du chi-deux}
% ────────────────────────────────────────────────────────────
\begin{frame}{4.3 — Distance du $\chi^2$}
\begin{definition}[Distance du $\chi^2$ entre profils lignes]
\[
d_{\chi^2}^2(\mathbf{L}_i, \mathbf{L}_{i'})
= \sum_{j=1}^J \frac{1}{f_{\bullet j}}
  \left(\frac{f_{ij}}{f_{i\bullet}} - \frac{f_{i'j}}{f_{i'\bullet}}\right)^2.
\]
C'est une distance euclidienne dans $\R^J$ avec \alert{métrique
$\mathbf{M} = \mathbf{D}_{1/f_{\bullet j}}$} (inverses des fréquences marginales).
\end{definition}

\begin{prop}[Propriétés essentielles]
\begin{enumerate}
  \item \textbf{Pondération} : les colonnes \emph{rares} pèsent plus lourd
        ($1/f_{\bullet j}$).
  \item \textbf{Équivalence distributionnelle} : si deux colonnes sont
        proportionnelles, leur fusion ne change pas les distances entre lignes.
  \item \textbf{Symétrie} : on définit de même la distance entre profils
        colonnes avec $\mathbf{M} = \mathbf{D}_{1/f_{i\bullet}}$.
\end{enumerate}
\end{prop}

\begin{rema}
\small La métrique du $\chi^2$ \alert{équilibre l'influence} des modalités peu
fréquentes mais informatives.
\end{rema}
\end{frame}

% ────────────────────────────────────────────────────────────
\subsection{4.4 Lien avec le test d'indépendance}
% ────────────────────────────────────────────────────────────
\begin{frame}{4.4 — Test d'indépendance et inertie totale}
\textbf{Hypothèse d'indépendance} : $f_{ij} = f_{i\bullet} \, f_{\bullet j}$.

\begin{theorem}[Statistique du $\chi^2$ et inertie]
La statistique de Pearson est :
\[
\chi^2 = n \sum_{i,j} \frac{(f_{ij} - f_{i\bullet} f_{\bullet j})^2}{f_{i\bullet} f_{\bullet j}}
       = n \, \Phi^2.
\]
Le coefficient $\Phi^2 = \chi^2 / n$ est exactement l'\alert{inertie totale} du
nuage des profils lignes (ou colonnes), pondérés par leurs marges :
\[
\boxed{\;
I_{\text{totale}}
= \sum_{i=1}^I f_{i\bullet}\, d_{\chi^2}^2(\mathbf{L}_i, \mathbf{L}_{\bullet})
= \Phi^2 = \chi^2 / n.
\;}
\]
\end{theorem}

\begin{rema}[Conséquence]
\small L'AFC \emph{décompose} la statistique du $\chi^2$ : chaque axe factoriel
explique une part de l'écart à l'indépendance.
\[
\Phi^2 = \sum_{k=1}^{r} \lambda_k, \qquad r = \min(I-1, J-1).
\]
\end{rema}
\end{frame}

% ────────────────────────────────────────────────────────────
\subsection{4.5 Formulation et résolution}
% ────────────────────────────────────────────────────────────
\begin{frame}{4.5 — Formulation comme SVD généralisée}
\textbf{Tableau centré} : $\widetilde{f}_{ij} = f_{ij} - f_{i\bullet} f_{\bullet j}$.

\begin{theorem}[Diagonalisation]
Soit $\mathbf{S}$ la matrice de terme général $\widetilde{f}_{ij}/\sqrt{f_{i\bullet} f_{\bullet j}}$.
La SVD $\mathbf{S} = \mathbf{U}\mathbf{\Sigma}\mathbf{V}^\top$ fournit les axes principaux
de l'AFC :
\begin{itemize}
  \item Valeurs propres : $\lambda_k = \sigma_k^2$ (vérifient $\lambda_k \leq 1$)
  \item Inertie totale : $\Phi^2 = \sum_k \lambda_k$
  \item Axes des lignes : $\mathbf{D}_{1/\sqrt{f_{i\bullet}}} \mathbf{U}$
  \item Axes des colonnes : $\mathbf{D}_{1/\sqrt{f_{\bullet j}}} \mathbf{V}$
\end{itemize}
\end{theorem}

\begin{block}{Cadre unificateur}
\small AFC $=$ ACP du triplet $(\mathbf{X}, \mathbf{M}, \mathbf{D}_p)$ avec :
\begin{itemize}\small
  \item $\mathbf{X}$ = tableau des profils lignes
  \item $\mathbf{M} = \mathbf{D}_{1/f_{\bullet j}}$ (métrique $\chi^2$)
  \item $\mathbf{D}_p = \mathbf{D}_{f_{i\bullet}}$ (pondération par marges)
\end{itemize}
\end{block}
\end{frame}

% ────────────────────────────────────────────────────────────
\subsection{4.6 Coordonnées factorielles}
% ────────────────────────────────────────────────────────────
\begin{frame}{4.6 — Coordonnées factorielles}
\begin{definition}[Coordonnées des lignes et des colonnes]
Sur le $k$-ième axe, les coordonnées sont :
\[
F_k(i) \quad\text{(coordonnée de la ligne } i\text{)}, \qquad
G_k(j) \quad\text{(coordonnée de la colonne } j\text{)}.
\]
\end{definition}

\begin{theorem}[Inertie expliquée par axe]
\[
\sum_{i=1}^I f_{i\bullet}\, F_k(i)^2 = \lambda_k,
\qquad
\sum_{j=1}^J f_{\bullet j}\, G_k(j)^2 = \lambda_k.
\]
Lignes et colonnes contribuent \emph{ensemble} à la même inertie $\lambda_k$.
\end{theorem}

\begin{rema}[Double point de vue]
\small L'AFC traite \alert{symétriquement} lignes et colonnes : on diagonalise
deux matrices semblables qui partagent les mêmes valeurs propres non nulles
(dualité). C'est ce qui justifie la \textbf{représentation simultanée}.
\end{rema}
\end{frame}

% ────────────────────────────────────────────────────────────
\subsection{4.7 Formules de transition}
% ────────────────────────────────────────────────────────────
\begin{frame}{4.7 — Formules barycentriques (transition)}
\begin{theorem}[Relations de transition]
Pour tout axe $k$ et avec $\lambda_k$ la valeur propre associée :
\[
\boxed{\;
F_k(i) = \frac{1}{\sqrt{\lambda_k}}
         \sum_{j=1}^J \frac{f_{ij}}{f_{i\bullet}}\, G_k(j),
\qquad
G_k(j) = \frac{1}{\sqrt{\lambda_k}}
         \sum_{i=1}^I \frac{f_{ij}}{f_{\bullet j}}\, F_k(i).
\;}
\]
\end{theorem}

\begin{block}{Lecture barycentrique}
\small La position d'une ligne $i$ est, à un facteur près, le \alert{barycentre}
des positions des colonnes $j$ pondérées par son profil $f_{ij}/f_{i\bullet}$.
Et inversement.
\end{block}

\begin{rema}[Interprétation graphique]
\small Une modalité ligne et une modalité colonne \alert{proches} sur le plan
factoriel sont dans une \textbf{relation d'attraction} : leur fréquence
conjointe dépasse l'attendu sous indépendance.
\end{rema}
\end{frame}

% ────────────────────────────────────────────────────────────
\subsection{4.8 Représentation simultanée}
% ────────────────────────────────────────────────────────────
\begin{frame}{4.8 — Représentation simultanée}
\begin{columns}[T]
\column{0.49\textwidth}
\textbf{Biplot AFC}\\
\small Sur le \emph{même} plan, on représente :
\begin{itemize}\small
  \item les modalités \alert{lignes} ($I$ points)
  \item les modalités \alert{colonnes} ($J$ points)
\end{itemize}

\textbf{Règles de lecture}
\begin{itemize}\small
  \item Deux \emph{lignes} proches : profils similaires
  \item Deux \emph{colonnes} proches : profils similaires
  \item Ligne $\times$ colonne proche : association forte
  \item Origine $=$ profil moyen
        (proche $=$ proche de l'indépendance)
\end{itemize}
\column{0.49\textwidth}
\begin{tikzpicture}[scale=0.95]
\draw[->, gray!50] (-2.5,0) -- (2.5,0) node[right, font=\tiny]{$F_1$};
\draw[->, gray!50] (0,-2) -- (0,2) node[above, font=\tiny]{$F_2$};

\fill[deepblue]   ( 1.7, 1.5)  circle(2.5pt) node[above, font=\tiny] {Cadre};
\fill[deepblue]   ( 0.4, 0.6)  circle(2.5pt) node[above, font=\tiny] {ProfInt};
\fill[deepblue]   (-1.0,-0.5)  circle(2.5pt) node[below, font=\tiny] {Employé};
\fill[deepblue]   (-2.0,-1.4)  circle(2.5pt) node[below, font=\tiny] {Ouvrier};

\fill[deeporange] (-2.0,-1.0)  circle(2.5pt) node[above, font=\tiny] {Bac};
\fill[deeporange] (-0.3, 0.4)  circle(2.5pt) node[right, font=\tiny] {Lic.};
\fill[deeporange] ( 1.2, 1.2)  circle(2.5pt) node[right, font=\tiny] {Mast.};
\fill[deeporange] ( 1.9, 1.6)  circle(2.5pt) node[right, font=\tiny] {Doct.};

\node[font=\tiny, gray, align=center] at (-2.0,1.7) {Lignes (CSP)\\Colonnes (Diplôme)};
\end{tikzpicture}
\end{columns}
\end{frame}

% ────────────────────────────────────────────────────────────
\subsection{4.9 Aides à l'interprétation}
% ────────────────────────────────────────────────────────────
\begin{frame}{4.9 — Aides à l'interprétation}
\begin{definition}[Qualité de représentation et contributions]
Pour la ligne $i$ sur l'axe $k$ :
\[
\cos^2_k(i) = \frac{F_k(i)^2}{\sum_{l \geq 1} F_l(i)^2},
\qquad
\text{ctr}_k(i) = \frac{f_{i\bullet}\, F_k(i)^2}{\lambda_k}.
\]
Définitions analogues pour les colonnes.
\end{definition}

\begin{block}{Règles de lecture}
\begin{itemize}\small
  \item \alert{Contribution élevée} : la modalité \emph{construit} l'axe
        (seuil pratique : ctr $> 1/I$ ou $1/J$)
  \item \alert{$\cos^2$ élevé} : la modalité est \emph{bien représentée}
  \item Une modalité avec $\cos^2$ faible doit être interprétée avec prudence
\end{itemize}
\end{block}

\begin{rema}
\small On peut décomposer le $\chi^2$ en sommes par cellule (\textbf{contribution
au $\chi^2$}) :
$\chi^2_{ij} = n \cdot (f_{ij} - f_{i\bullet}f_{\bullet j})^2 / (f_{i\bullet}f_{\bullet j})$,
ce qui complète l'AFC.
\end{rema}
\end{frame}

% ────────────────────────────────────────────────────────────
\subsection{4.10 Reconstitution du tableau}
% ────────────────────────────────────────────────────────────
\begin{frame}{4.10 — Reconstitution et approximation}
\begin{theorem}[Formule de reconstitution]
Le tableau de fréquences peut être reconstruit à partir des $r$ axes
principaux :
\[
f_{ij} = f_{i\bullet} f_{\bullet j}
        \left(1 + \sum_{k=1}^r \frac{F_k(i)\, G_k(j)}{\sqrt{\lambda_k}}\right).
\]
\end{theorem}

\begin{block}{Lecture}
\begin{itemize}\small
  \item Le terme $f_{i\bullet} f_{\bullet j}$ est l'\alert{indépendance théorique}.
  \item La somme corrective décompose les écarts à l'indépendance par axe.
  \item Tronquer la somme à $q < r$ axes donne la \textbf{meilleure
        approximation de rang $q$} (Eckart-Young).
\end{itemize}
\end{block}

\begin{rema}[Comparaison à l'ACP]
\small AFC et ACP partagent la même structure : SVD généralisée + Eckart-Young.
La différence est uniquement dans le \emph{triplet} $(\mathbf{X}, \mathbf{M}, \mathbf{D}_p)$.
\end{rema}
\end{frame}

% ────────────────────────────────────────────────────────────
\subsection{4.11 Modalités supplémentaires}
% ────────────────────────────────────────────────────────────
\begin{frame}{4.11 — Modalités et tableaux supplémentaires}
\textbf{Modalité supplémentaire} : projetée sur les axes \emph{après} construction,
sans contribuer à leur calcul.

\begin{block}{Formules de projection}
Pour une ligne supplémentaire $i^*$ de profil $\mathbf{L}_{i^*}$ :
\[
F_k(i^*) = \frac{1}{\sqrt{\lambda_k}}
         \sum_{j=1}^J \frac{f_{i^*j}}{f_{i^*\bullet}}\, G_k(j).
\]
Symétriquement pour une colonne supplémentaire.
\end{block}

\begin{columns}[T]
\column{0.49\textwidth}
\textbf{Cas d'usage}
\begin{itemize}\small
  \item Comparer un échantillon test à une population de référence
  \item Suivre une modalité dans le temps (cohorte)
  \item Caractériser un axe par une variable externe
\end{itemize}
\column{0.49\textwidth}
\begin{alertblock}{Attention}
\small Une modalité supplémentaire \emph{rare} ($f \to 0$) peut être très
mal positionnée : on l'interprète avec son $\cos^2$.
\end{alertblock}
\end{columns}
\end{frame}

% ────────────────────────────────────────────────────────────
\subsection{4.12 Code R}
% ────────────────────────────────────────────────────────────
\begin{frame}[fragile]{4.12 — AFC avec FactoMineR (et ca)}
\begin{lstlisting}
library(FactoMineR); library(factoextra)

# Tableau de contingence (CSP x Diplome)
tab <- as.table(rbind(
  Cadre   = c(Bac=12,  Lic=65,  Mast=180, Doct=95),
  ProfInt = c(Bac=88,  Lic=110, Mast=50,  Doct=5),
  Employe = c(Bac=230, Lic=45,  Mast=12,  Doct=1),
  Ouvrier = c(Bac=195, Lic=8,   Mast=2,   Doct=0)))

# Test d'independance du chi2
chisq.test(tab)

# Analyse Factorielle des Correspondances
res <- CA(tab, graph = FALSE)

# Valeurs propres et inertie
get_eigenvalue(res)               # lambda_k, % d'inertie
fviz_eig(res, addlabels = TRUE)   # scree plot

# Representation simultanee (biplot)
fviz_ca_biplot(res, repel = TRUE)

# cos2 et contributions
res$row$cos2; res$col$cos2
res$row$contrib; res$col$contrib

# Decomposition du chi2 par cellule
chisq.test(tab)$residuals^2       # ~ contributions au chi2
\end{lstlisting}
\end{frame}

% ────────────────────────────────────────────────────────────
\subsection{4.13 Récapitulatif}
% ────────────────────────────────────────────────────────────
\begin{frame}{4.13 — Récapitulatif du chapitre}
\begin{columns}[T]
\column{0.49\textwidth}
\textbf{À retenir}
\begin{itemize}\small
  \item AFC $=$ ACP des \alert{profils} avec métrique $\chi^2$
  \item Inertie totale $\Phi^2 = \chi^2 / n$
  \item Lignes et colonnes : mêmes valeurs propres
  \item Représentation \alert{simultanée} sur un biplot
  \item Formules de transition $\Rightarrow$ lecture barycentrique
  \item Reconstitution comme \emph{somme} d'écarts à l'indépendance
\end{itemize}

\column{0.49\textwidth}
\begin{block}{Garanties théoriques}
\begin{itemize}\small
  \item SVD généralisée : optimalité (Eckart-Young)
  \item Conservation de l'inertie totale
  \item Décomposition orthogonale du $\chi^2$
\end{itemize}
\end{block}

\begin{alertblock}{Domaine d'application}
\small AFC $\to$ \alert{2 variables qualitatives}.
Pour $p > 2$ variables qualitatives, on utilise l'\textbf{ACM}
(Analyse des Correspondances Multiples), chap. 5.
\end{alertblock}
\end{columns}
\end{frame}

%% Chapitre — Analyse des Correspondances Multiples (ACM)
%% À compléter
\section{Analyse des Correspondances Multiples (ACM)}

%% Chapitre — Classification
%% À compléter
\section{Classification}

%% Chapitre — Analyse Factorielle Discriminante (AFD)
%% À compléter
\section{Analyse Factorielle Discriminante (AFD)}

% ============================================================
%  Chapitre 8 — Méthodes avancées et applications
% ============================================================
\section{Méthodes Avancées et Applications}

% ────────────────────────────────────────────────────────────
\subsection{8.1 Vue d'ensemble}
% ────────────────────────────────────────────────────────────
\begin{frame}{8.1 — Au-delà des méthodes de base}
\begin{columns}[T]
\column{0.55\textwidth}
\textbf{Méthodes complémentaires}
\begin{itemize}\small
  \item \alert{MDS} — partir de \emph{distances} et reconstituer un nuage
  \item \alert{Analyse canonique} — relier deux \emph{groupes} de variables
  \item \alert{AFM / AFMD} — analyser plusieurs \emph{tableaux} simultanément
  \item \alert{PLS} — régression sur variables très corrélées
  \item \alert{Variantes non linéaires} — Kernel PCA, Isomap, t-SNE
\end{itemize}

\vspace{0.2cm}
\textbf{Cadre} : toutes ces méthodes étendent la machinerie spectrale
(Ch. 2) à des situations \alert{multi-tableaux} ou \alert{non linéaires}.
\column{0.42\textwidth}
\begin{tikzpicture}[scale=0.9, font=\scriptsize,
  box/.style={draw=deepblue, fill=deepblue!10, rounded corners,
              minimum width=2.4cm, minimum height=0.55cm, align=center}]
\node[box] (acp) at (0,3) {ACP};
\node[box] (afc) at (0,2) {AFC};
\node[box] (acm) at (0,1) {ACM};

\node[box, fill=deeporange!10, draw=deeporange] (mds)  at (4,3) {MDS};
\node[box, fill=deeporange!10, draw=deeporange] (cca)  at (4,2) {CCA};
\node[box, fill=deeporange!10, draw=deeporange] (afm)  at (4,1) {AFM};
\node[box, fill=deeporange!10, draw=deeporange] (pls)  at (4,0) {PLS};

\draw[->, gray] (acp) -- (mds);
\draw[->, gray] (acp) -- (cca);
\draw[->, gray] (afc) -- (afm);
\draw[->, gray] (acp) -- (pls);

\node[font=\tiny, gray, align=center] at (-1.2, 2.0) {méthodes\\de base};
\node[font=\tiny, gray, align=center] at (5.5, 1.5) {extensions};
\end{tikzpicture}
\end{columns}
\end{frame}

% ────────────────────────────────────────────────────────────
\subsection{8.2 MDS classique}
% ────────────────────────────────────────────────────────────
\begin{frame}{8.2 — Multidimensional Scaling (MDS) classique}
\begin{definition}[Problème du MDS]
Étant donnée une matrice de \textbf{distances} (ou dissimilarités)
$\mathbf{D} = (d_{ij}) \in \R^{n \times n}$ entre $n$ objets, trouver
des coordonnées $\mathbf{Y} \in \R^{n \times q}$ telles que :
\[
\|\mathbf{y}_i - \mathbf{y}_j\| \approx d_{ij}, \quad \forall i, j.
\]
\end{definition}

\begin{theorem}[MDS classique de Torgerson]
Soit $\mathbf{B} = -\frac{1}{2}\,\mathbf{H}\mathbf{D}^{(2)}\mathbf{H}$ où
$\mathbf{D}^{(2)} = (d_{ij}^2)$ et $\mathbf{H} = \mathbf{I}_n - \frac{1}{n}\mathbf{1}_n \mathbf{1}_n^\top$
est la matrice de centrage. Si $\mathbf{B}$ est positive, ses $q$ premiers
vecteurs propres fournissent $\mathbf{Y} = \mathbf{U}_q \mathbf{\Lambda}_q^{1/2}$ qui
restitue exactement les $\|\mathbf{y}_i - \mathbf{y}_j\| = d_{ij}$.
\end{theorem}

\begin{remark}[Lien MDS / ACP]
\small Si $d_{ij}$ est la distance euclidienne entre lignes d'un tableau $\mathbf{X}$,
le MDS classique \alert{coïncide} avec l'ACP non normée (à une rotation près).
\end{remark}
\end{frame}

% ────────────────────────────────────────────────────────────
\subsection{8.3 MDS non métrique}
% ────────────────────────────────────────────────────────────
\begin{frame}{8.3 — MDS non métrique (Shepard-Kruskal)}
\textbf{Idée} : ne préserver que l'\alert{ordre} des dissimilarités, pas leurs
valeurs numériques.

\begin{definition}[Stress de Kruskal]
On minimise
\[
\text{Stress}(\mathbf{Y}) = \sqrt{\frac{\sum_{i < j}(\hat{d}_{ij} - \|\mathbf{y}_i - \mathbf{y}_j\|)^2}{\sum_{i < j} \|\mathbf{y}_i - \mathbf{y}_j\|^2}}
\]
où $\hat{d}_{ij}$ sont des \emph{disparities} respectant l'ordre des $d_{ij}$
(régression isotone).
\end{definition}

\begin{block}{Cas d'usage}
\begin{itemize}\small
  \item Dissimilarités issues d'un \emph{questionnaire} (rangs)
  \item Données \emph{phylogénétiques}, \emph{psychométriques}
  \item Toute situation où la métrique est ordinale
\end{itemize}
\end{block}

\begin{remark}[Critère de qualité]
\small Stress $< 0.05$ : excellent ; $< 0.10$ : bon ; $> 0.20$ : médiocre.
\end{remark}
\end{frame}

% ────────────────────────────────────────────────────────────
\subsection{8.4 Analyse canonique}
% ────────────────────────────────────────────────────────────
\begin{frame}{8.4 — Analyse Canonique des Corrélations (CCA)}
\textbf{Cadre} : deux groupes de variables sur les mêmes individus,
$\mathbf{X} \in \R^{n \times p}$ et $\mathbf{Y} \in \R^{n \times q}$.

\begin{definition}[Variables canoniques]
On cherche des combinaisons linéaires
\[
U = \mathbf{X}\mathbf{a}, \qquad V = \mathbf{Y}\mathbf{b}
\]
maximisant la \alert{corrélation} :
\[
\rho_1 = \max_{\mathbf{a}, \mathbf{b}}\; \Cor(U, V).
\]
\end{definition}

\begin{theorem}[Solution]
Les vecteurs $\mathbf{a}, \mathbf{b}$ sont solutions du système couplé
\[
\mathbf{V}_{XX}^{-1}\mathbf{V}_{XY}\mathbf{V}_{YY}^{-1}\mathbf{V}_{YX}\, \mathbf{a} = \rho^2 \mathbf{a},
\quad
\mathbf{V}_{YY}^{-1}\mathbf{V}_{YX}\mathbf{V}_{XX}^{-1}\mathbf{V}_{XY}\, \mathbf{b} = \rho^2 \mathbf{b}.
\]
On obtient $\min(p, q)$ couples canoniques $(\mathbf{a}_k, \mathbf{b}_k)$ deux à deux orthogonaux.
\end{theorem}

\begin{remark}
\small La CCA est l'analogue \emph{symétrique} de la régression linéaire : elle
relie deux groupes \emph{sans hiérarchie}.
\end{remark}
\end{frame}

% ────────────────────────────────────────────────────────────
\subsection{8.5 Analyse Factorielle Multiple (AFM)}
% ────────────────────────────────────────────────────────────
\begin{frame}{8.5 — Analyse Factorielle Multiple (AFM)}
\textbf{Cadre} : un même ensemble d'individus décrit par
\alert{plusieurs groupes} de variables, possiblement hétérogènes
(quantitatives + qualitatives).

\begin{definition}[Pondération AFM]
Soit $\mathbf{X} = [\mathbf{X}^{(1)} | \mathbf{X}^{(2)} | \cdots | \mathbf{X}^{(J)}]$
la concaténation de $J$ blocs. L'AFM applique une ACP \alert{pondérée}
au tableau global, avec un poids $1/\lambda_1^{(j)}$ par bloc, où
$\lambda_1^{(j)}$ est la première valeur propre de l'ACP du bloc $j$.
\end{definition}

\begin{theorem}[Équilibrage]
Avec cette pondération, l'inertie maximale d'un axe contribuée par chaque bloc
est \alert{normalisée à 1}. Aucun bloc ne domine la première composante
en raison de sa taille ou de son échelle.
\end{theorem}

\begin{block}{Sorties spécifiques}
\begin{itemize}\small
  \item \textbf{Coordonnées} des individus, des variables, et des \emph{groupes}
  \item \textbf{RV-coefficient} : mesure d'accord entre 2 groupes
  \item \textbf{Lg-coefficient} : généralisation au cas multi-groupes
\end{itemize}
\end{block}
\end{frame}

% ────────────────────────────────────────────────────────────
\subsection{8.6 Régression PLS}
% ────────────────────────────────────────────────────────────
\begin{frame}{8.6 — Régression PLS (Partial Least Squares)}
\textbf{Problème} : régresser $\mathbf{Y}$ sur $\mathbf{X}$ avec $p \gg n$ ou
$\mathbf{X}$ très corrélé. La régression OLS échoue ($\mathbf{X}^\top \mathbf{X}$ singulière).

\begin{definition}[Composantes PLS]
PLS construit successivement des composantes $\mathbf{t}_k = \mathbf{X}\mathbf{w}_k$
maximisant la \alert{covariance} (et non la corrélation) avec $\mathbf{Y}$ :
\[
\mathbf{w}_k = \argmax_{\|\mathbf{w}\|=1}\; \Cov^2(\mathbf{X}\mathbf{w}, \mathbf{Y}),
\]
puis on régresse $\mathbf{Y}$ sur $\mathbf{t}_1, \ldots, \mathbf{t}_q$.
\end{definition}

\begin{block}{Différence avec ACP + régression}
\begin{itemize}\small
  \item ACP : maximise $\Var(\mathbf{X}\mathbf{w})$ — \alert{ignore} $\mathbf{Y}$
  \item PLS : maximise $\Cov(\mathbf{X}\mathbf{w}, \mathbf{Y})$ — combine variance et corrélation
  \item Avantage en \textbf{prédiction}
\end{itemize}
\end{block}

\begin{remark}
\small La PLS est centrale en \emph{chimiométrie} (spectres infrarouge $\to$
concentration), \emph{génomique} (microarrays) et plus généralement quand
$p \gg n$.
\end{remark}
\end{frame}

% ────────────────────────────────────────────────────────────
\subsection{8.7 Méthodes non linéaires}
% ────────────────────────────────────────────────────────────
\begin{frame}{8.7 — Variantes non linéaires}
\begin{columns}[T]
\column{0.49\textwidth}
\textbf{Kernel PCA (Schölkopf, 1998)}
\begin{itemize}\small
  \item ACP dans un espace de Hilbert via un noyau $K(x,y)$
  \item Capte des structures \emph{courbes}
  \item Diagonalise la matrice $\mathbf{K} = (K(x_i, x_j))$
\end{itemize}

\textbf{Isomap}
\begin{itemize}\small
  \item MDS classique sur la distance \emph{géodésique}
  \item Préserve la \emph{topologie} d'une variété
\end{itemize}

\textbf{LLE (Locally Linear Embedding)}
\begin{itemize}\small
  \item Reconstruit chaque point comme combinaison linéaire de ses voisins
  \item Pas de noyau, complexité $O(n^2)$
\end{itemize}
\column{0.49\textwidth}
\textbf{t-SNE et UMAP}
\begin{itemize}\small
  \item Préservent les \emph{voisinages} probabilistes
  \item Utiles pour la \emph{visualisation 2D}
  \item \alert{Pas factoriels} (pas de coordonnées globales)
\end{itemize}

\begin{block}{Critère de choix}
\begin{itemize}\small
  \item Données linéaires $\to$ ACP / AFC
  \item Variétés courbes connues $\to$ Kernel PCA
  \item Visualisation seulement $\to$ t-SNE / UMAP
  \item Spectres ou prédiction $\to$ PLS
\end{itemize}
\end{block}
\end{columns}
\end{frame}

% ────────────────────────────────────────────────────────────
\subsection{8.8 Études de cas}
% ────────────────────────────────────────────────────────────
\begin{frame}{8.8 — Études de cas métier}
\begin{columns}[T]
\column{0.49\textwidth}
\textbf{1. Marketing — segmentation client}
\begin{itemize}\small
  \item Variables RFM + qualitatives
  \item AFDM ou ACM + HCPC
  \item Profils de classes $\to$ campagnes ciblées
\end{itemize}

\textbf{2. Finance — risque de crédit}
\begin{itemize}\small
  \item Variables continues (revenus, ratios)
  \item AFD : LDA pour score
  \item Validation : courbe ROC, KS
\end{itemize}

\textbf{3. Sociologie — enquête}
\begin{itemize}\small
  \item Tableau de questionnaire
  \item ACM + variables sociales suppl.
  \item Analyse des positions sociales (Bourdieu)
\end{itemize}
\column{0.49\textwidth}
\textbf{4. Sensoriel — analyse sensorielle}
\begin{itemize}\small
  \item AFM sur jurys × descripteurs
  \item Cartes de produits
  \item Convergence des panels
\end{itemize}

\textbf{5. Bio-informatique — micro-arrays}
\begin{itemize}\small
  \item $p \gg n$, gènes corrélés
  \item PLS ou Sparse PCA
  \item Identification de signatures
\end{itemize}

\textbf{6. NLP — corpus textuel}
\begin{itemize}\small
  \item Matrice termes $\times$ documents
  \item AFC / LSA
  \item Topics latents
\end{itemize}
\end{columns}
\end{frame}

% ────────────────────────────────────────────────────────────
\subsection{8.9 Limites et extensions modernes}
% ────────────────────────────────────────────────────────────
\begin{frame}{8.9 — Limites et extensions modernes}
\textbf{Limites des méthodes factorielles classiques}
\begin{itemize}\small
  \item Hypothèse de \alert{linéarité} (sauf Kernel PCA)
  \item Sensibilité aux \alert{outliers} (variantes robustes)
  \item Sélection des \alert{axes} encore largement subjective
  \item Mauvaise adaptation aux données \alert{très volumineuses} ($n \gg 10^6$)
\end{itemize}

\vspace{0.2cm}
\textbf{Extensions modernes}
\begin{itemize}\small
  \item \textbf{Sparse PCA} : axes parcimonieux et interprétables
  \item \textbf{Robust PCA} (Candès et al.) : décompose $\mathbf{X} = \mathbf{L} + \mathbf{S}$
        (low-rank + sparse)
  \item \textbf{Functional PCA} : pour données fonctionnelles ($y(t)$)
  \item \textbf{Tensor decomposition} : Tucker, CP, pour tableaux 3D+
  \item \textbf{Auto-encodeurs} (réseaux de neurones) : ACP non linéaire profonde
  \item \textbf{Topological Data Analysis} (TDA) : persistance, mapper
\end{itemize}

\begin{block}{Filiation moderne}
\small Les méthodes factorielles restent \alert{le squelette} interprétable que
les approches récentes (deep learning) cherchent à intégrer ou à remplacer.
\end{block}
\end{frame}

% ────────────────────────────────────────────────────────────
\subsection{8.10 Code R}
% ────────────────────────────────────────────────────────────
\begin{frame}[fragile]{8.10 — Méthodes avancées en R}
\begin{lstlisting}
# 1) MDS classique
data(eurodist)             # distances villes
mds <- cmdscale(eurodist, k = 2, eig = TRUE)
plot(mds$points, type = "n"); text(mds$points, labels = names(eurodist))

# 2) MDS non metrique
library(MASS)
mds.nm <- isoMDS(eurodist, k = 2)

# 3) Analyse canonique (CCA)
library(CCA)
res.cca <- cc(X1, X2)
res.cca$cor                # correlations canoniques

# 4) AFM (Multiple Factor Analysis)
library(FactoMineR)
res.afm <- MFA(don, group = c(2, 5, 3), type = c("c","c","n"),
               name.group = c("ratios","prix","categorie"))
plot(res.afm, choix = "ind")
plot(res.afm, choix = "group")

# 5) Regression PLS
library(pls)
res.pls <- plsr(y ~ ., data = train, ncomp = 10, validation = "CV")
summary(res.pls); plot(RMSEP(res.pls))

# 6) Kernel PCA
library(kernlab)
res.kpca <- kpca(~ ., data = X, kernel = "rbfdot")
plot(rotated(res.kpca)[, 1:2])
\end{lstlisting}
\end{frame}

% ────────────────────────────────────────────────────────────
\subsection{8.11 Récapitulatif}
% ────────────────────────────────────────────────────────────
\begin{frame}{8.11 — Récapitulatif du chapitre}
\begin{columns}[T]
\column{0.49\textwidth}
\textbf{Méthodes vues}
\begin{itemize}\small
  \item \textbf{MDS} (classique, non métrique)
  \item \textbf{Analyse canonique} (CCA)
  \item \textbf{AFM} (multi-tableaux)
  \item \textbf{PLS} (régression $p \gg n$)
  \item \textbf{Kernel PCA, Isomap, LLE} (non linéaire)
  \item \textbf{Sparse / Robust PCA}, fonctionnel, tensoriel
\end{itemize}
\column{0.49\textwidth}
\begin{block}{Ligne directrice}
\begin{itemize}\small
  \item Toutes ces méthodes \emph{généralisent} la SVD
  \item Elles répondent à des \emph{questions différentes}
  \item Choisir selon : type de données, structure recherchée, objectif
\end{itemize}
\end{block}

\begin{alertblock}{Vers la pratique}
\small La maîtrise de ces extensions repose sur : (1) compréhension du
chapitre 2, (2) choix de la \emph{métrique} et de la \emph{pondération},
(3) validation par les TPs.
\end{alertblock}
\end{columns}
\end{frame}

% ────────────────────────────────────────────────────────────
\subsection{8.12 Conclusion générale du cours}
% ────────────────────────────────────────────────────────────
\begin{frame}{8.12 — Conclusion générale du cours}
\begin{columns}[T]
\column{0.49\textwidth}
\textbf{Ce que vous savez maintenant}
\begin{enumerate}\small
  \item Décomposer un tableau de données en termes d'\alert{inertie}
  \item Choisir la \emph{bonne} méthode factorielle
        selon le type de données
  \item Démontrer les théorèmes fondamentaux
        (Spectral, SVD, Eckart-Young, Huygens)
  \item Lire un plan factoriel et un cercle des corrélations
  \item Construire une typologie validée
        (HCPC + caractérisation)
  \item Mettre en œuvre tout cela en \textbf{R}
        (FactoMineR, factoextra, MASS)
\end{enumerate}
\column{0.49\textwidth}
\begin{block}{Trois principes à retenir}
\begin{enumerate}\small
  \item \textbf{Géométrie d'abord} : tout est inertie, projection, distance.
  \item \textbf{Une seule machine} : SVD généralisée du triplet
        $(\mathbf{X}, \mathbf{M}, \mathbf{D}_p)$.
  \item \textbf{Toujours interpréter} : les chiffres ne parlent pas seuls,
        les variables et le métier oui.
\end{enumerate}
\end{block}

\begin{alertblock}{Vers la suite}
\small Les TPs en R consolident la maîtrise pratique. Le polycopié
fournit les démonstrations détaillées. Bonne fin de cours !
\end{alertblock}
\end{columns}
\end{frame}

%% Bibliographie
%% À compléter
\section*{Bibliographie}


\end{document}
